Single-window vulnerability prioritization studies evaluate methods at one point in time, but enterprise operations teams patch continuously across rolling maintenance windows. As high-priority CVEs are remediated, the scoring landscape changes: exploit-likelihood signals such as the Exploit Prediction Scoring System (EPSS) degrade as their top-ranked CVEs disappear from the patch backlog, while host-hygiene signals may remain discriminative longer because hygiene posture changes more slowly than CVE remediation state. This paper examines whether \textit{HygienePrio}'s Precision@K advantage over EPSS-only persists across six consecutive bi-weekly maintenance windows in a pre-registered multi-window simulation on the Extended Enterprise Hygiene Dataset Artifact (EEHDA) synthetic fleet.

The central finding is that EPSS-only prioritization decays from mean P@50 $= 0.331$ at Window~1 to $0.034$ at Window~6, an 89.6\% relative drop, while HygienePrio-full maintains mean P@50 $\geq 0.499$ at every window. The absolute P@50 advantage of HygienePrio-full over EPSS-only grows from $26.4$ percentage points at Window~1 to $46.5$ percentage points at Window~6. HygienePrio-full outperforms EPSS-only in all $150$ of $150$ window-seed pairs (25 evaluation seeds $\times$ 6 windows). By Window~2, HRS-only ($0.229$) has overtaken EPSS-only ($0.149$), and the ordering does not reverse for the remainder of the simulation. Hygiene Risk Score signals are temporally resilient because patch debt, Active Directory exposure, and telemetry freshness evolve on longer timescales than per-CVE exploit likelihood, making them suitable primary signals for multi-window scheduling. All results are bounded to the synthetic evaluation context; external validation on real fleet data with longitudinal exploitation outcome data is required before deployment.
